% TOP_PROBLEM: {"id":30007632,"problem_number":"LOCAL-30007632","title":"Measuring employer wage-setting power","category_id":21,"category":{"id":21,"name":"economics","display_name":"Economics","description":"Open economic theory statements, published research agendas, and proposed scoped specifications; question provenance and historical priority are qualified per record.","slug":"economics","order_index":21,"created_at":"2026-10-08T00:00:00Z"},"status":"research_question","external_url":null,"legacy_ids":[],"published":false,"statement":"In the explicitly specified setting: How much of the wage–productivity gap reflects employer market power rather than amenities, measurement error, worker sorting, or incentive contracts? The mathematical statement above fixes the target, assumptions and scope of this catalogue formulation.","economics":{"original_id":121,"field":"Labor markets and work","kind":"identification","formalization_basis":"adapted_research_question","source_status":"research_agenda","priority_status":"earliest_found","priority_note":"Earliest directly inspected qualifying passage in this audit; earlier origins were not established. A review or research-agenda statement does not imply original priority.","earliest_verified_reference":{"authors":["Patrick Kline"],"title":"Labor Market Monopsony: Fundamentals and Frontiers","year":2025,"url":"https://eml.berkeley.edu/~pkline/papers/HoLE_vol6_monopsony.pdf","doi":"10.1016/bs.heslab.2025.07.007","locator":"Introduction, printed pp. 658–660; section 3.1.2, printed p. 684","evidence_summary":"Identifying wage-setting power requires restrictions beyond productivity-to-wage passthrough; several empirical frontiers remain."},"current_reference":{"authors":["Patrick Kline"],"title":"Labor Market Monopsony: Fundamentals and Frontiers","year":2025,"url":"https://eml.berkeley.edu/~pkline/papers/HoLE_vol6_monopsony.pdf","doi":"10.1016/bs.heslab.2025.07.007","locator":"Introduction, printed pp. 658–660; section 3.1.2, printed p. 684","evidence_summary":"Identifying wage-setting power requires restrictions beyond productivity-to-wage passthrough; several empirical frontiers remain."},"formalization_note":"The cited passage establishes a research direction, not this exact mathematical statement. The notation, population, policy menu, assumptions, and estimands are an original adaptation for this catalogue; no universal unsolved theorem or source priority is claimed."},"statusQualification":"Published question or research agenda verified at the cited locator. The mathematical specification is adapted for this catalogue and includes additional assumptions; source publication does not establish first-ever priority or certify this exact model remains unsolved. The cited passage establishes a research direction, not this exact mathematical statement. The notation, population, policy menu, assumptions, and estimands are an original adaptation for this catalogue; no universal unsolved theorem or source priority is claimed.","statusReviewedAt":"2026-10-08"}
% FIRST_POSTED: 2026-10-08
% ENTRY_KIND: statement_only
% !TeX program = lualatex
\documentclass[11pt]{article}
\usepackage{fontspec}
\usepackage[a4paper,margin=25mm]{geometry}
\usepackage{amsmath,amssymb,mathtools}
\usepackage{xurl}
\usepackage[hidelinks]{hyperref}
\newcommand{\sref}[2]{\hyperlink{source:#1:#2}{[S#2]}}
\setlength{\emergencystretch}{3em}
\begin{document}
% problemId:30007632
\section*{Measuring employer wage-setting power}\label{30007632}
\subsection{Definitions and mathematical statement}
Fix a labor market, a time interval, and a population of employer--worker matches. Let \(w_{ij}>0\) be worker \(i\)'s hourly compensation at firm \(j\), and let \(p_{ij}>0\) be the marginal revenue product of that worker's labor, holding the firm's other inputs fixed. Define the observed wedge \(d_{ij}=\log p_{ij}-\log w_{ij}\). Let \(a_{ij}\) denote job amenities, \(s_i\) worker skills, and \(c_{ij}\) contract incentives. In a structural model \(M\), let \(w^{C}_{ij}(M)>0\) be compensation when the firm's ability to restrict labor supply or bargain over rents is removed while preferences, production, and aggregate resources remain fixed. Define it as a feasible offer to the same worker when a baseline match disappears; report cases with no such offer separately. The market-power component is \(m_{ij}(M)=\log w^{C}_{ij}(M)-\log w_{ij}\). Given matched records, production information, and specified exogenous shocks, identify or sharply bound \(E[m_{ij}(M)]\) over models consistent with the data. Establish which restrictions distinguish this component from amenity compensation, sorting, productivity measurement error, and incentive pay. Report sensitivity to the definition of competition and to equilibrium reallocation. The task concerns the identified set in a declared market; neither a wage wedge nor a passthrough coefficient alone constitutes the answer.

\par\textbf{Formulation and scope.} The cited passage establishes a research direction, not this exact mathematical statement. The notation, population, policy menu, assumptions, and estimands are an original adaptation for this catalogue; no universal unsolved theorem or source priority is claimed.

\par\textbf{Open-question provenance.} An explicit question or research agenda was verified at the source locator. This does not by itself certify that every later version remains unresolved \sref{30007632}{1}.

\par\textbf{Historical priority.} Earliest directly inspected qualifying passage in this audit; earlier origins were not established. A review or research-agenda statement does not imply original priority.

\subsection{Short English statement}
In the explicitly specified setting: How much of the wage–productivity gap reflects employer market power rather than amenities, measurement error, worker sorting, or incentive contracts? The mathematical statement above fixes the target, assumptions and scope of this catalogue formulation.

\subsection{Sources}
\begin{itemize}\raggedright
\item[\textnormal{[S1]}]\hypertarget{source:30007632:1}{} Patrick Kline. 2025. Labor Market Monopsony: Fundamentals and Frontiers. Introduction, printed pp. 658–660; section 3.1.2, printed p. 684.
\url{https://eml.berkeley.edu/~pkline/papers/HoLE_vol6_monopsony.pdf}.
DOI: 10.1016/bs.heslab.2025.07.007.
{\small\textit{Earliest verified question/agenda reference; historical priority qualified below. Identifying wage-setting power requires restrictions beyond productivity-to-wage passthrough; several empirical frontiers remain.}}
\end{itemize}
\end{document}
